\section{Human Validation}
\label{sec:human}

NLI-based faithfulness scoring captures entailment between answer claims
and retrieved context, but reviewers reasonably question whether the
metric tracks the kind of faithfulness humans actually care about. This
section reports a Prolific study designed to (a) calibrate NLI scores
against human judgement and (b) test whether the coherence paradox
survives when faithfulness is judged by humans rather than a model.

\subsection{Study design}

We sample 500 (case, condition) items stratified across the five
head-to-head conditions (§\ref{sec:headtohead}) — 100 items per
condition, balanced within each condition across the six datasets in
proportion to the multidataset query mix. Each item is double-annotated
by independent Prolific workers (no overlap), yielding 1{,}000 annotation
tasks. The full study budget is approximately £450 at the standard UK
Prolific rate (£9/hour, 1.5~min per item).

Annotators see the question, the retrieved passages, and the candidate
answer; they do \emph{not} see the condition label or any system
metadata. The rubric (full text in
\texttt{results/human\_validation/study\_v1/annotation\_guide.md}) asks
for:

\begin{itemize}
  \item A binary \textbf{Faithful} judgement (every claim supported by
        the passages, yes/no), with explicit handling of honest
        abstention (``I cannot find this in the passages'' = Faithful=YES
        when correct).
  \item A 1--5 \textbf{Correctness} rating.
  \item A 1--5 \textbf{Helpfulness} rating.
  \item An optional one-sentence comment.
\end{itemize}

Workers who fail a 3-item calibration set with $<$80\% agreement against
the gold rubric are rejected and the items re-issued.

\subsection{Preregistered reporting plan}
\label{sec:human_plan}

The human-validation outputs will be reported in three blocks, each
preregistered before the Prolific launch. First, inter-annotator
agreement: Fleiss' $\kappa$ on the binary \textbf{Faithful} judgement
and on the two ordinal 1--5 ratings (\textbf{Correctness},
\textbf{Helpfulness}), interpreted under the conventional bands
($<$0.20 poor, 0.21--0.40 fair, 0.41--0.60 moderate, 0.61--0.80
substantial, $>$0.81 almost perfect). Second, NLI-vs-human calibration:
Spearman and Pearson correlations between the NLI faithfulness score
and each human dimension, with $p$-values on $\rho$ and sample sizes
restricted to items with at least two valid annotations and a defined
NLI score. Third, per-condition human-faithfulness rates side-by-side
with NLI-faithfulness rates for \texttt{baseline}, \texttt{hcpc\_v1},
\texttt{hcpc\_v2}, \texttt{crag}, and \texttt{selfrag}.

\subsection{Does the paradox survive human evaluation?}

The load-bearing question is whether HCPC-v1 — which raises NLI-measured
faithfulness above baseline on parts of the §\ref{sec:paradox} table
while \emph{lowering} answer quality — also looks worse to humans, or
whether the paradox is an artifact of the NLI scorer. The
per-condition reporting block (human-faith rate, NLI-faith rate, and
mean correctness for each of the five conditions) is the one the
preregistered plan above reads directly on this question. The paradox
is considered to survive human evaluation if \texttt{hcpc\_v1}'s human
faithfulness rate is lower than \texttt{baseline}'s on the same matched
query set; otherwise the §\ref{sec:paradox} effect is treated as an NLI
artifact.

\subsection{Reporting commitments}

We commit in advance to reporting the human-evaluation outcome regardless
of direction:

\begin{itemize}
  \item If $\rho_{\text{NLI},\text{human}} \geq 0.5$ \emph{and} the
        paradox replicates (HCPC-v1 human-faith $<$ baseline human-faith),
        the §\ref{sec:paradox} story stands.
  \item If $\rho_{\text{NLI},\text{human}} < 0.3$, we narrow the
        faithfulness claims to ``NLI-measured faithfulness'' throughout
        and treat the human study as evidence that NLI is the wrong
        oracle.
  \item If the paradox does \emph{not} replicate under humans (HCPC-v1
        human-faith $\geq$ baseline human-faith), we report it as a
        falsification: the §\ref{sec:paradox} effect was an NLI
        artifact and HCPC-v2's recovery would need a different
        justification.
\end{itemize}

The annotation pipeline, rubric, and analysis script are released as
part of ContextCoherenceBench (§\ref{sec:benchmark}) so the study can be
replicated or extended on additional conditions.
